\documentclass[12pt]{article}

\usepackage{amsmath, amssymb}
\usepackage{geometry}
\usepackage{setspace}
\usepackage{graphicx}
\usepackage{hyperref}

\geometry{margin=1in}
\setstretch{1.2}

\title{CSE400 -- Fundamentals of Probability in Computing\\
Lecture 16 Scribe\\
Joint Random Variables and Joint Distributions}
\author{}
\date{}

\begin{document}

\maketitle

\section{Introduction and Motivation}

In many scientific and engineering problems, we must analyze \textbf{two or more random variables simultaneously} rather than studying each variable independently. This motivates the concept of \textbf{joint probability}, which allows us to compute probabilities of events involving multiple variables occurring together.

Examples where multiple variables are studied together include:

\begin{itemize}
\item Height and weight of giraffes
\item IQ and birthweight of children
\item Frequency of exercise and rate of heart disease
\item Air pollution level and respiratory illness rate
\item Number of Facebook friends and age of members
\end{itemize}

These examples demonstrate that two variables may have a relationship, meaning the value of one variable may provide information about the other.

\section{Applications of Joint Probability}

\subsection{Smart Transportation}

Consider two random variables:

\[
X = \text{Vehicle Speed}
\]

\[
Y = \text{Traffic Density}
\]

A key question is whether vehicle speed can be high when traffic density is also high. This illustrates that speed and density may be \textbf{correlated random variables}.

Applications include:

\begin{itemize}
\item Adaptive traffic signals
\item Congestion prediction
\item Autonomous vehicle routing
\end{itemize}

\subsection{Wireless Network Performance}

In wireless communication systems we consider:

\[
X = \text{Signal-to-Noise Ratio (SNR)}
\]

\[
Y = \text{Data Throughput}
\]

Throughput depends on SNR but may also be affected by network congestion and interference. Thus even when SNR is high, throughput may not always be high. This indicates a \textbf{conditional relationship between variables}.

Applications include:

\begin{itemize}
\item 5G scheduling
\item Adaptive modulation
\item Network planning
\end{itemize}

\section{Concept of Joint Distribution}

When two random variables are studied together, they possess a \textbf{joint distribution}. This distribution allows us to:

\begin{itemize}
\item Compute probabilities involving multiple variables
\item Understand relationships between variables
\item Analyze dependence structures
\end{itemize}

If the variables are independent, the analysis becomes simpler. When they are not independent, dependence can be measured using quantities such as covariance and correlation.

\section{Discrete Case: Joint Probability Mass Function}

\subsection{Setup}

Suppose $X$ and $Y$ are discrete random variables:

\[
X = \{x_1, x_2, \dots , x_n\}
\]

\[
Y = \{y_1, y_2, \dots , y_m\}
\]

The ordered pair $(X,Y)$ takes values in the \textbf{product set}

\[
\{(x_1,y_1),(x_1,y_2),\dots,(x_n,y_m)\}.
\]

\subsection{Definition}

The \textbf{joint probability mass function (PMF)} is defined as

\[
p(x_i,y_j) = P(X=x_i, Y=y_j)
\]

This represents the probability that $X$ takes value $x_i$ and $Y$ simultaneously takes value $y_j$.

\section{Joint PMF Table Representation}

The joint PMF can be organized into a two-dimensional probability table where:

\begin{itemize}
\item Rows correspond to values of $X$
\item Columns correspond to values of $Y$
\item Each entry contains the probability $p(x_i,y_j)$
\end{itemize}

\subsection{Key Properties}

\textbf{Property 1:}

\[
0 \leq p(x_i,y_j) \leq 1
\]

\textbf{Property 2:}

\[
\sum_{i=1}^{n} \sum_{j=1}^{m} p(x_i,y_j) = 1
\]

These conditions ensure that the table represents a valid probability distribution.

\section{Example: Rolling Two Dice}

\subsection{Experiment}

Roll two fair dice.

Define random variables:

\[
X = \text{value on the first die}
\]

\[
T = \text{total sum of both dice}
\]

There are $36$ equally likely outcomes. Thus each ordered pair outcome $(i,j)$ has probability

\[
P((i,j)) = \frac{1}{36}.
\]

\subsection{Observations}

From the joint probability table we observe:

\begin{itemize}
\item Entries are either $0$ or $\frac{1}{36}$
\item Some combinations are impossible and therefore have probability $0$
\end{itemize}

If $X=1$, possible totals are

\[
T = 2,3,4,5,6,7
\]

Each corresponding outcome occurs with probability $\frac{1}{36}$.

\subsection{Dependence}

From the joint table we observe that

\[
X \text{ and } T \text{ are not independent}
\]

because knowing $X$ restricts the possible values of $T$.

\section{Continuous Case: Joint Probability Density Function}

The continuous case generalizes the discrete case.

\begin{center}
\begin{tabular}{c|c}
Discrete Case & Continuous Case \\
\hline
Discrete values & Continuous intervals \\
Joint PMF & Joint PDF \\
Summations & Double integrals \\
\end{tabular}
\end{center}

If

\[
X \in [a,b], \quad Y \in [c,d]
\]

then $(X,Y)$ lies in the rectangular region

\[
[a,b] \times [c,d].
\]

\subsection{Definition}

A function $f(x,y)$ is called the \textbf{joint probability density function} if

\[
P((X,Y) \text{ lies in a small rectangle}) \approx f(x,y)\,dx\,dy
\]

\subsection{Key Properties}

\textbf{Non-negativity}

\[
f(x,y) \ge 0
\]

\textbf{Normalization}

\[
\int_c^d \int_a^b f(x,y)\,dx\,dy = 1
\]

Note that $f(x,y)$ may be greater than $1$ because it represents a density, not a probability.

\section{Worked Continuous Example}

Given

\[
f(x,y) = 4xy
\]

for

\[
0 \le x \le 1, \quad 0 \le y \le 1
\]

Thus $(X,Y)$ lies in the unit square $[0,1]^2$.

\subsection{Step 1: Check Validity}

\[
\int_0^1 \int_0^1 4xy \, dx\,dy
\]

First integrate with respect to $x$:

\[
\int_0^1 4xy\,dx = 2y
\]

Then integrate with respect to $y$:

\[
\int_0^1 2y\,dy = [y^2]_0^1 = 1
\]

Thus the function is a valid joint PDF.

\subsection{Step 2: Compute Probability}

Let

\[
A = \{X < 0.5, Y > 0.5\}
\]

Then

\[
P(A) = \int_0^{0.5} \int_{0.5}^{1} 4xy\,dy\,dx
\]

Evaluating the integral gives

\[
P(A) = \frac{3}{16}.
\]

\section{Conclusion}

Key takeaways from this lecture include:

\begin{itemize}
\item Joint Random Variable models events involving multiple random variables.
\item The joint PMF describes discrete joint distributions.
\item The joint PDF describes continuous joint distributions.
\item Joint distributions allow analysis of relationships between variables.
\item Dependence between variables can be studied through joint distributions.
\end{itemize}

\end{document}
